% !PRINTED 339
\weberpage{339}
%<<CONT
\gesperrt{$\alpha$ ist, so erh\"alt man ein engeres Intervall $(\alpha', \beta')$
von denselben Eigenschaften, wenn man}
\weq{4}{
\alpha' = \alpha - \frac{f(\alpha)(\beta - \alpha)}{f(\beta) - f(\alpha)}, \quad \beta' = \beta - \frac{f(\beta)}{f'(\beta)}
}
\gesperrt{setzt.}

Um diese Resultate der geometrischen Anschauung analytisch
 zu beweisen, nehmen wir an, es sei im Intervall $f'(x)$ und
$f''(x)$ positiv, also $\alpha < \beta$, und
\[
f(\alpha) < 0, \quad f(\beta) > 0.
\]

Wir beschr\"anken uns auf die Betrachtung dieses einen
Falles, da die drei anderen genau in derselben Weise zu behandeln
 sind.

Wir bilden die Function
\[
\varphi(x) = \frac{f(\beta) - f(x)}{\beta - x}
\]
und deren erste Derivirte
\[
\varphi'(x) = \frac{-(\beta - x)f'(x) + f(\beta) - f(x)}{(\beta - x)^2}.
\]
Der Z\"ahler dieses Ausdruckes
\[
\psi(x) = -(\beta - x)f'(x) + f(\beta) - f(x)
\]
verschwindet f\"ur $x = \beta$, und seine Derivirte ist
\[
\psi'(x) = -(\beta - x)f''(x),
\]
also im Intervall negativ. Daraus folgt, dass $\psi(x)$ im Intervall
mit wachsendem $x$ abnimmt und daher, weil es f\"ur den gr\"ossten
Werth $x = \beta$ verschwindet, positiv bleibt. Folglich ist auch
$\varphi'(x)$ positiv und $\varphi(x)$ w\"achst im Intervall mit wachsendem $x$
best\"andig. Wir haben demnach
\weq{5}{
\begin{gathered}
\frac{f(\beta) - f(\alpha)}{\beta - \alpha} < \frac{f(\beta) - f(x)}{\beta - x} < f'(\beta) \\
\alpha < x < \beta.
\end{gathered}
}
Setzen wir nun
\[
\alpha' = \alpha - \frac{f(\alpha)(\beta - \alpha)}{f(\beta) - f(\alpha)}, \quad \beta' = \beta - \frac{f(\beta)}{f'(\beta)},
\]
so wird
\weq{6}{
\beta' - \alpha' = f(\beta)\left\{\frac{\beta - \alpha}{f(\beta) - f(\alpha)} - \frac{1}{f'(\beta)}\right\},
}
und nach (5) folgt, dass diese Differenz positiv ist, also
\[
\alpha < \alpha' < \beta' < \beta.
\]
% 22*
